% ---------------------------------------------------------
\subsection*{Axiom of Completeness}
% ---------------------------------------------------------


\begin{remark*}[Why completeness is needed]
The ordered field axioms alone do not prevent ``holes'' in the number line.
The rationals $\mathbb{Q}$ satisfy all field and order axioms, yet fail
completeness. Consider the set
\[
S = \{x \in \mathbb{Q} : x^2 < 2\}.
\]
This set is nonempty (e.g.\ $1 \in S$) and bounded above in $\mathbb{Q}$
(e.g.\ $2$ is an upper bound). Yet $\sup S$ does not exist as a rational
number: the candidate $\sqrt{2}$ is irrational. The set $S$ has no least
upper bound in $\mathbb{Q}$ --- a hole sits exactly where $\sqrt{2}$ should be.

The Completeness Axiom asserts that $\mathbb{R}$ has no such holes:
every nonempty bounded-above set has a supremum \emph{inside} $\mathbb{R}$.
This single axiom is what distinguishes $\mathbb{R}$ from $\mathbb{Q}$,
and it underlies every major theorem in analysis.
\end{remark*}


\begin{definitionbox}{Definition (Least Upper Bound Property)}
\begin{definition}[Least Upper Bound Property]
\label{def:least-upper-bound-property}
Let $(S,\le)$ be an ordered set. The set $S$ has the \emph{Least Upper Bound Property} if and only if every nonempty subset of $S$ that is bounded above in $S$ has a supremum in $S$.
\end{definition}
\end{definitionbox}

\begin{remark*}[Standard quantified statement]
\[
\begin{aligned}
&\text{Let }(S,\le)\text{ be an ordered set}.\\
&S\text{ has the Least Upper Bound Property}\\
&\Longleftrightarrow
\forall A\subseteq S\;
\Bigl[
\bigl(A\neq\varnothing \land \exists u\in S\;\forall x\in A\;(x\le u)\bigr)
\rightarrow
\exists s\in S\;\operatorname{Supremum}(s,A)
\Bigr].
\end{aligned}
\]
\end{remark*}

\begin{remark*}[Definition predicate reading]
\[
\operatorname{LeastUpperBoundProperty}(S)\coloneqq
\forall A\subseteq S\;
\Bigl[
\bigl(A\neq\varnothing \land \exists u\in S\;\forall x\in A\;(x\le u)\bigr)
\rightarrow
\exists s\in S\;\operatorname{Supremum}(s,A)
\Bigr].
\]
\end{remark*}

\begin{remark*}[Negated quantified statement]
\[
\begin{aligned}
&\text{Let }(S,\le)\text{ be an ordered set}.\\
&S\text{ does not have the Least Upper Bound Property}\\
&\Longleftrightarrow
\exists A\subseteq S\;
\Bigl[
A\neq\varnothing
\land
\exists u\in S\;\forall x\in A\;(x\le u)
\land
\forall s\in S\;\neg\operatorname{Supremum}(s,A)
\Bigr].
\end{aligned}
\]
\end{remark*}

\begin{remark*}[Negation predicate reading]
\[
\neg\operatorname{LeastUpperBoundProperty}(S)\Longleftrightarrow
\exists A\subseteq S\;
\Bigl[
A\neq\varnothing
\land
\exists u\in S\;\forall x\in A\;(x\le u)
\land
\forall s\in S\;\neg\operatorname{Supremum}(s,A)
\Bigr].
\]
\end{remark*}

\begin{remark*}[Failure modes]
The property fails when there is a nonempty subset of the ambient ordered set that has at least one upper bound but has no least upper bound inside that same ambient set.
\end{remark*}

\begin{remark*}[Failure mode decomposition]
\[
\neg\operatorname{LeastUpperBoundProperty}(S)
=
\underbrace{\exists A\subseteq S\;(A\neq\varnothing)}_{\text{nonempty witness set}}
\;\land\;
\underbrace{\exists u\in S\;\forall x\in A\;(x\le u)}_{\text{bounded above in }S}
\;\land\;
\underbrace{\forall s\in S\;\neg\operatorname{Supremum}(s,A)}_{\text{no supremum in }S}.
\]
\end{remark*}

\begin{remark*}[Interpretation]
The Least Upper Bound Property says that upper-bounded nonempty subsets do not point toward missing ceiling elements. Whenever the order supplies at least one upper bound, it also supplies a smallest upper bound in the same ambient set. Failure therefore indicates a hole in the ordered structure rather than a defect in the subset alone.
\end{remark*}

\begin{dependencies}
\begin{itemize}
  \item \hyperref[ax:completeness-of-reals]{Completeness of $\mathbb{R}$}
\end{itemize}
\end{dependencies}






\begin{axiombox}{Axiom (Completeness of $\mathbb{R}$)}
\begin{axiom}[Completeness of $\mathbb{R}$]
\label{ax:completeness-of-reals}
The ordered field $\mathbb{R}$ has the Least Upper Bound Property.
\end{axiom}
\end{axiombox}

\begin{remark*}[Standard quantified statement]
\[
\operatorname{LeastUpperBoundProperty}(\mathbb{R}).
\]
Equivalently,
\[
\begin{aligned}
&\forall A\subseteq\mathbb{R}\;
\Bigl[
\bigl(A\neq\varnothing \land \exists u\in\mathbb{R}\;\forall x\in A\;(x\le u)\bigr)
\rightarrow
\exists s\in\mathbb{R}\;\operatorname{Supremum}(s,A)
\Bigr].
\end{aligned}
\]
\end{remark*}

\begin{remark*}[Axiom predicate reading]
\[
\operatorname{LeastUpperBoundProperty}(\mathbb{R}).
\]
\end{remark*}

\begin{remark*}[Negated quantified statement]
\[
\begin{aligned}
&\neg\operatorname{LeastUpperBoundProperty}(\mathbb{R})\\
&\Longleftrightarrow
\exists A\subseteq\mathbb{R}\;
\Bigl[
A\neq\varnothing
\land
\exists u\in\mathbb{R}\;\forall x\in A\;(x\le u)
\land
\forall s\in\mathbb{R}\;\neg\operatorname{Supremum}(s,A)
\Bigr].
\end{aligned}
\]
\end{remark*}

\begin{remark*}[Negation predicate reading]
\[
\neg\operatorname{LeastUpperBoundProperty}(\mathbb{R})\Longleftrightarrow
\exists A\subseteq\mathbb{R}\;
\Bigl[
A\neq\varnothing
\land
\exists u\in\mathbb{R}\;\forall x\in A\;(x\le u)
\land
\forall s\in\mathbb{R}\;\neg\operatorname{Supremum}(s,A)
\Bigr].
\]
\end{remark*}

\begin{remark*}[Failure modes]
Completeness would fail if some nonempty subset of $\mathbb{R}$ were bounded above in $\mathbb{R}$ but had no supremum in $\mathbb{R}$.
\end{remark*}

\begin{remark*}[Failure mode decomposition]
\[
\neg\operatorname{LeastUpperBoundProperty}(\mathbb{R})
=
\underbrace{\exists A\subseteq\mathbb{R}\;(A\neq\varnothing)}_{\text{nonempty witness set}}
\;\land\;
\underbrace{\exists u\in\mathbb{R}\;\forall x\in A\;(x\le u)}_{\text{bounded above in }\mathbb{R}}
\;\land\;
\underbrace{\forall s\in\mathbb{R}\;\neg\operatorname{Supremum}(s,A)}_{\text{missing real supremum}}.
\]
\end{remark*}

\begin{remark*}[Interpretation]
The completeness axiom asserts that the real line has no order gaps of the least-upper-bound kind. Nonempty sets of real numbers that are bounded above always determine a real number sitting exactly at the least possible ceiling. This is the structural property that distinguishes $\mathbb{R}$ from ordered fields such as $\mathbb{Q}$.
\end{remark*}

\NoLocalDependencies





